\wordchapter{5}{Results}

\wordsection{5.1}{Pre-Pilot Trajectories}

Figure 6 compares changes before and after the pilots relative to the last pre-pilot year. Each point averages six crop-specific contrasts, each calculated against its eligible comparison crops. Pre-pilot contrasts near zero would be consistent with parallel observed changes. A pattern in which those contrasts move toward or away from zero indicates differences in observed growth before insurance became available. The reference-year contrast equals zero by construction.

Pre-pilot yield contrasts range from approximately −0.04 to 0.02 log points. The joint p-value is 0.92. This pattern provides less evidence of divergent observed growth than the area and production series, although a failure to reject with six treated crops does not establish parallel trends.

Area and production show larger differences in pre-pilot growth. Contrasts of approximately −0.2 log points in earlier years approach zero near the reference year, indicating that treated crops were gaining relative to comparison crops before insurance became available. The exploratory joint p-values are 0.084 for area and 0.0535 for production. The production value is close to 0.05. The observed trajectories are therefore more informative than a sharp distinction between rejection and non-rejection. Their similar patterns are consistent with the accounting relationship among area, yield, and production.

\begin{wordfigure}{6}{Event-Time Contrasts for Cultivated Area, Yield, and Production}
\wordimage{5.39583in}{figures/fig6_event_time_contrasts.png}
\wordnotes{\emph{Notes:} Event-time comparisons are relative to the last pre-pilot year (open circle), and are averaged over the 6 treated crops; comparison crops are not including the other treated crops. Bars are exploratory 95 percent intervals from 9,999 Webb multiplier draws; no nominal coverage is determined for this design (Appendix A). The dotted line running through the middle is the pre-pilot and post-pilot years and the transition years between the pilot and expansion.}
\end{wordfigure}

This pre-pilot divergence also affects the choice of baseline. A reference year close to introduction measures subsequent changes from a later point in the relative gains of treated crops. Section 5.3 examines how earlier reference periods change the contrasts.

\wordsection{5.2}{Aggregate and Crop-Level Contrasts}

Table 4 presents the main contrasts using two comparison pools. The first includes the full pool. The second excludes spring napa cabbage and spring radish, whose survey method changed in 2012.

\begin{wordtable}{4}{Mean Post-Period Contrasts}
\begin{wordtabular}{L{0.422}C{0.193}C{0.193}C{0.193}}
\wtop
\hdl{0.422}{} & \hd{0.193}{(1)\par Area} & \hd{0.193}{(2)\par Yield} & \hd{0.193}{(3)\par Production} \\
\whead
\panelgap
\textit{Panel A. Full comparison pool} & & & \\
Mean contrast & 0.011 & −0.012 & −0.002 \\
 & (0.141) & (0.056) & (0.168) \\
 & [−0.253, 0.275] & [−0.118, 0.093] & [−0.318, 0.315] \\
Comparison series & 22 & 22 & 22 \\
\panelgap
\textit{Panel B. Excluding two spring series} & & & \\
Mean contrast & −0.069 & −0.014 & −0.084 \\
 & (0.139) & (0.059) & (0.169) \\
 & [−0.332, 0.193] & [−0.125, 0.096] & [−0.406, 0.238] \\
Comparison series & 20 & 20 & 20 \\
\wbottom
\end{wordtabular}

\wordnotes{\emph{Notes}: Log outcomes; six treated crops; ten target years. In contrast to Figure 6, target years start with the year of expansion. Panel B does not include radish or napa cabbage, which are common in spring. Score-based standard errors (in parentheses) and 95 percent intervals (in brackets) from 9,999 Webb multiplier draws; their nominal coverage is not determined for this design (Appendix A). * p < 0.10, ** p < 0.05, *** p < 0.01.}
\end{wordtable}

Using the full pool, I obtain mean contrasts of 0.011 log points for area, −0.012 for yield, and −0.002 for production, with respective exploratory p-values of 0.94, 0.84, and 0.99. Excluding the two spring series changes the contrasts to −0.069, −0.014, and −0.084. Area reverses sign and production moves further below zero, while yield changes little. The production interval under the full pool, [−0.318, 0.315] log points, includes both substantial decreases and substantial increases.

Figure 7 reports the crop-specific means underlying the aggregate contrasts.

\begin{wordfigure}{7}{Mean Contrasts by Treated Crop}
\wordimage{6.2in}{figures/fig7_crop_contrasts.png}
\wordnotes{\emph{Notes}: Each point averages ten crop-specific contrasts. Circles are based on the full comparison pool and triangles on the comparison pool without spring napa cabbage and spring radish. Log points.}
\end{wordfigure}

The small aggregate combines large differences across crops. Production contrasts are 0.406 for onion, 0.154 for sweet potato, and −0.634 for red pepper. Under equal weights, these crops contribute approximately 0.068, 0.026, and −0.106 log points, respectively. Dropping one crop at a time moves the aggregate between −0.083 and 0.125 (Appendix Table A1). Individual crops thus strongly influence the mean. The comparison alone does not distinguish insurance responses from other shocks or differences in untreated growth.

\wordsection{5.3}{Sensitivity to Comparison Choices}

Sensitivities with respect to reference period, comparison pool and timing convention are reported in Table 5 in the same format as in Table 4.

\begin{wordtable}{5}{Sensitivity to the Reference Period, Comparison Crops, and Timing}
\begin{wordtabular}[1.30]{L{0.423}C{0.188}C{0.188}C{0.198}}
\wtop
\hdl{0.423}{} & \hd{0.188}{(1)\par Area} & \hd{0.188}{(2)\par Yield} & \hd{0.198}{(3)\par Production} \\
\whead
\panelgap
\textit{Panel A. Reference period: mean of last five pre-pilot years} & & & \\
Contrast & 0.051 & 0.008 & 0.059 \\
 & (0.154) & (0.049) & (0.170) \\
 & [−0.244, 0.346] & [−0.085, 0.102] & [−0.264, 0.382] \\
\panelgap
\textit{Panel B. Reference period: mean of all pre-pilot years} & & & \\
Contrast & 0.238 & −0.029 & 0.209 \\
 & (0.202) & (0.076) & (0.208) \\
 & [−0.144, 0.621] & [−0.172, 0.114] & [−0.184, 0.601] \\
\panelgap
\textit{Panel C. Comparison crops: field crops only} & & & \\
Contrast & −0.082 & −0.011 & −0.094 \\
 & (0.198) & (0.083) & (0.250) \\
 & [−0.450, 0.285] & [−0.166, 0.143] & [−0.557, 0.369] \\
\panelgap
\textit{Panel D. Comparison crops: excluding selected vegetable series} & & & \\
Contrast & −0.113 & −0.005 & −0.118 \\
 & (0.160) & (0.067) & (0.196) \\
 & [−0.414, 0.187] & [−0.130, 0.120] & [−0.487, 0.250] \\
\panelgap
\textit{Panel E. Timing: product-table years} & & & \\
Contrast & 0.015 & −0.012 & 0.002 \\
 & (0.139) & (0.056) & (0.165) \\
 & [−0.247, 0.277] & [−0.116, 0.092] & [−0.308, 0.313] \\
\panelgap
\textit{Panel F. Timing: target years begin with pilot} & & & \\
Contrast & 0.001 & −0.010 & −0.009 \\
 & (0.117) & (0.050) & (0.136) \\
 & [−0.220, 0.221] & [−0.102, 0.082] & [−0.267, 0.249] \\
\wbottom
\end{wordtabular}

\wordnotes{\emph{Notes:} Log outcomes; 10 target years; full comparison pool unless noted; 6 treated crops. The standard errors are given in parentheses and exploratory 95 percent intervals are given in brackets in the score-based estimates from 9,999 Webb multiplier draws (Appendix A). * p < 0.10, ** p < 0.05, *** p < 0.01. For multi-year specifications, the comparison crops must be observed and unexposed for each year that is referenced (Table 4, Panel A). Panel D does not include napa cabbage and radish forms, carrot or green onions. The product-table calendar years take the place of harvest-year timing in panel E. While untreated, other crops can be used as comparison crops; the ten target years begin with the pilot year in panel F.}
\end{wordtable}

All intervals in Table 5 include zero. The smallest exploratory p-value is 0.244, for area when all pre-pilot years form the reference. Point estimates nevertheless vary substantially. Earlier reference periods raise the production contrast to 0.059 and 0.209, consistent with the relative gains before introduction shown in Figure 6. Restricting the comparison pool lowers it to −0.094 or −0.118. The two alternative timing conventions give production contrasts of 0.002 and −0.009. Standard errors for area and production range from 0.12 to 0.25. The comparisons therefore remain imprecise across these specifications.

Pre-pilot area weights also change the aggregate (Table 6). The production contrast falls from −0.002 to −0.105 with the full comparison pool and from −0.084 to −0.186 after excluding the two spring series. Soybean and red pepper account for 65 percent of these weights. Onion and sweet potato, which have positive contrasts, receive less weight. None of the weighted contrasts is statistically significant. This weighting provides a different summary of the same crop contrasts; it does not estimate the percentage change in national production.

\begin{wordtable}{6}{Contrasts under Pre-Pilot Area Weights}
\begin{wordtabular}[1.225]{L{0.425}C{0.195}C{0.195}C{0.195}}
\wtop
\hdl{0.425}{} & \hd{0.195}{(1)\newline Area} & \hd{0.195}{(2)\newline Yield} & \hd{0.195}{(3)\newline Production} \\
\whead
\panelgap
\textit{Panel A. Full comparison pool} & & & \\
Contrast & −0.108 & 0.004 & −0.105 \\
 & (0.140) & (0.092) & (0.189) \\
 & [−0.368, 0.152] & [−0.158, 0.165] & [−0.439, 0.229] \\
\panelgap
\textit{Panel B. Excluding two spring series} & & & \\
Contrast & −0.186 & 0.001 & −0.186 \\
 & (0.135) & (0.095) & (0.192) \\
 & [−0.433, 0.060] & [−0.168, 0.169] & [−0.527, 0.156] \\
\wbottom
\end{wordtabular}

\wordnotes{\emph{Notes}: Rows and symbols as in Table 5. Contrasts with equal weights are reported in Table 4. Area weights are soybean 39.0, red pepper 25.7, garlic 13.4, sweet potato 7.6, corn 7.2, and onion 6.9 percent; equal weights are 16.7 percent. Comparison crops remain equally weighted within each contrast.}
\end{wordtable}

\wordsection{5.4}{Municipal Rice Comparison}

Table 7 compares the rice cohorts over 2000-2008. Early entrants had a mean planted area of 14,425 hectares, compared with 13,666 hectares among later entrants. Mean slopes of log area were approximately −0.011 per year in both cohorts, although area levels were more dispersed among early entrants. Welch tests do not reject equality of the reported means at the 10 percent level.

\begin{wordtable}{7}{Pre-Pilot Characteristics of Early and Later Rice Municipalities, 2000-2008}
\begin{wordtabular}{L{0.496}C{0.180}C{0.179}C{0.145}}
\wtop
\hdl{0.496}{} & \hd{0.180}{(1)\par Early entrants} & \hd{0.179}{(2)\par Later entrants} & \hd{0.145}{(3)\par Difference} \\
\whead
Mean planted area, 2000-2008 (ha) & 14,425 & 13,666 & 759 \\
 & (5,871) & (1,998) & (1,572) \\
Annual growth of planted area, 2000-2008 & −0.011 & −0.011 & −0.000 \\
 & (0.011) & (0.009) & (0.004) \\
Change in log planted area, 2000-2008 & −0.079 & −0.075 & −0.003 \\
 & (0.077) & (0.069) & (0.030) \\
Dispersion around the pre-pilot trend & 0.017 & 0.012 & 0.005 \\
 & (0.012) & (0.003) & (0.003) \\
Municipalities & 17 & 9 &  \\
\wbottom
\end{wordtabular}

\wordnotes{\emph{Notes}: Columns (1) and (2) report means across municipalities, with standard deviations in parentheses. Column (3) reports the difference in means, with its Welch standard error in parentheses. Annual growth is the least-squares slope of log planted area on year; dispersion around the pre-pilot trend is the standard deviation of the residuals from that trend. * p < 0.10, ** p < 0.05, *** p < 0.01, from Welch t-tests. Data from Statistics Korea (n.d.).}
\end{wordtable}

The main estimate from equation (7) is −0.021 log points (Table 8). Its 95 percent interval, [−0.064, 0.022], corresponds to approximately −6.2 to +2.2 percent. Figure 8 reports pre-pilot coefficients between approximately −0.001 and 0.010, with a joint p-value of 0.52. Assigning a placebo period of 2006-2008 within the pre-pilot sample gives an estimate of 0.001 (p ≈ 0.96).

\begin{wordtable}{8}{Rice Planted Area in Pilot and Later-Entering Municipalities, 2000-2010}
\begin{wordtabular}{L{0.331}C{0.196}C{0.473}}
\wtop
\hdl{0.331}{} & \hd{0.196}{(1)\par Main sample} & \hd{0.473}{(2)\par Excluding municipalities adjoining weirs} \\
\whead
\textbf{Early entry × Post} & −0.021 & −0.002 \\
 & (0.022) & (0.021) \\
 & [−0.064, 0.022] & [−0.045, 0.039] \\
Municipalities (early / later) & 17 / 9 & 14 / 8 \\
Municipality-years & 286 & 242 \\
Municipality and year FE & Yes & Yes \\
\wbottom
\end{wordtabular}

\wordnotes{\emph{Notes}: The dependent variable is log paddy-rice planted area. Municipality-clustered (CR1) standard errors in parentheses. Intervals in brackets are 95 percent intervals, significance levels are calculated around the null-imposed wild cluster bootstrap-t with Webb weights, 9,999 draws and municipality clusters. * p < 0.10, ** p < 0.05, *** p < 0.01. Gumi, Sangju, Naju and Uiseong are municipalities neighbouring the weirs constructed under the Four Major Rivers Restoration Project (K-water, n.d.-a, n.d.-b) and were excluded from the numbers in column (2).}
\end{wordtable}

The estimate changes to −0.002 (p = 0.940) when I exclude Gumi, Sangju, Naju, and Uiseong, which adjoin weirs built under the Four Major Rivers Restoration Project. This result shows sensitivity to the municipalities included. Because the exclusion changes the sample, it does not separately identify the effect of river works. As an additional inference check, province-clustered CR1 standard errors yield p = 0.387 using a t distribution with seven degrees of freedom. Inference at this level relies on only eight province clusters.

\begin{wordfigure}{8}{Event-Study Coefficients for Rice Planted Area, 2000-2010}
\wordimage{6.2in}{figures/fig8_rice_event_study.png}
\wordnotes{\emph{Notes}: Fixed effect coefficients of log paddy-rice planted area based on the regression with municipality and year fixed effects and early-entry status-year interaction fixed effects (reference year is 2008, open circle). Bars represent 95 percent confidence intervals using t(25) critical values and municipality-clustered (CR1) standard errors; they are different than the wild bootstrap intervals in Table 8. The sample includes 17 municipalities with earlier entries and 9 with later entries and that have more than 2000.}
\end{wordfigure}

\wordsection{5.5}{Comparability of Treated and Comparison Crops}

Table 9 shows that treated and comparison crops followed different pre-pilot growth paths. Across the six treated and 22 contributing comparison crops, mean log-area slopes were −0.016 and −0.046 per year, respectively. Mean production slopes were −0.007 and −0.031. Yield slopes were closer, at 0.008 and 0.015. The standardized differences are 0.83 for area, 0.64 for production, and −0.37 for yield. The larger gaps for area and production reinforce the divergence shown in Figure 6.

\begin{wordtable}{9}{Pre-Pilot Growth of Treated and Comparison Crops, 1998-2007}
\begin{wordtabular}{L{0.352}C{0.170}C{0.170}C{0.145}C{0.163}}
\wtop
\hdl{0.352}{} & \hd{0.170}{(1)\par Treated crops} & \hd{0.170}{(2)\par Comparison crops} & \hd{0.145}{(3)\par Difference} & \hd{0.163}{(4)\par Std. difference} \\
\whead
Annual growth of cultivated area & −0.016 & −0.046 & 0.030 & 0.833 \\
 & (0.024) & (0.045) &  &  \\
Annual growth of yield & 0.008 & 0.015 & −0.007 & −0.371 \\
 & (0.021) & (0.014) &  &  \\
Annual growth of production & −0.007 & −0.031 & 0.023 & 0.643 \\
 & (0.023) & (0.046) &  &  \\
Crop series & 6 & 22 &  &  \\
\wbottom
\end{wordtabular}

\wordnotes{\emph{Notes}: Unweighted means and standard deviations (across crops) of crop-specific least-squares slopes of log outcomes on year over 1998-2007. The difference is the treated mean - comparison mean. Standardized difference is that difference divided by the square root of the average of the two group variances; it is descriptive, NOT a test statistic. There are eight annual observations of large and small green onions and ten of the other crops.}
\end{wordtable}

The national results are contingent on what history of growth is used for the comparison. Note that they are sensitive to the choice of reference period or comparison pool, as changing the reference period or comparison pool changes those histories. The municipal analysis filters out variations in crop type and is nevertheless based on the comparison of regional growth trajectories.
